% Targeted replacements for thesis PDF version (6).
% Copy each passage into the indicated location, retaining existing labels.
% This is a replacement pack, NOT a chapter to include wholesale.
% The complete version-(6) source project is not available locally.
% See audit.md for the equation audit, terminology changes, and open checks.

% ----------------------------------------------------------------------
% SECTION 3.4.1: append after the definition of SHAP importance.
% ----------------------------------------------------------------------
The SHAP signal describes the fitted post-hoc detector, not a causal effect
of a feature on data quality. Real records are labelled 1 and synthetic
records 0, so a positive attribution moves the explained output towards the
real class, while a negative attribution moves it towards the synthetic
class, relative to the explainer's reference. Importance uses the magnitude
of these attributions. For a categorical feature, signed one-hot
contributions are first summed within each record and the absolute value is
taken afterwards. The mean is calculated only over the selected correctly
identified synthetic holdout records. If this subset is empty, the
implementation returns a zero importance vector rather than an undefined mean.

% ----------------------------------------------------------------------
% SECTION 3.4.2: append after Eq. (3.4).
% ----------------------------------------------------------------------
If the audit-data interquartile range is zero or non-finite, the scale is
replaced by the population standard deviation of the same real audit values.
If that value is also zero or non-finite, the scale is set to one. The
categorical mismatch is the square-root Jensen--Shannon distance with
base-2 logarithms, not the Jensen--Shannon divergence.

% ----------------------------------------------------------------------
% SECTION 3.5: replace Eq. (3.8), keeping its original equation label.
% ----------------------------------------------------------------------
\begin{equation}
    \rho_j^{(a)}=
    \begin{cases}
        C_j^{(a)}/\sum_{k\in\mathcal{J}_a}C_k^{(a)},
        & j\in\mathcal{J}_a\text{ and }\sum_{k\in\mathcal{J}_a}C_k^{(a)}>0,\\
        0, & \text{otherwise}.
    \end{cases}
\end{equation}

% SECTION 3.5: replace the closing explanatory paragraph with this passage.
The contribution of SHAP has three stages. First, the importance signal
contributes to the combined feature scores and can therefore change which
features enter the top-$k$ set $\mathcal{J}_a$. Second, it affects their
relative priorities $\rho_j^{(a)}$. Third, those priorities determine how
strongly the distributional deficits of the selected features contribute to
each training row's score and hence its sampling probability. SHAP values
are not themselves row weights, are not added to the GAN loss, and do not
provide gradients to the generator. All weighted variants use signals
derived once from the common A0 audit, without iterative re-explanation.

In A3, $C_j^{(\mathrm{A3})}=\widehat I_j$: SHAP is the only signal used for
\emph{feature prioritisation}. This does not mean that the complete A3
weighting procedure is independent of distributional measurements. A3 still
uses the same one-sided region deficits $g_{jr}$ as A2, A4, and A5 to
identify underrepresented values of the selected features. A3 therefore
tests SHAP-only priorities within a shared deficit-based sampling framework.
It excludes the global mismatch and tail signals from feature scoring,
but it does not remove the region-deficit mechanism.

% ----------------------------------------------------------------------
% SECTION 3.7: optional clarification immediately after sampling equation.
% ----------------------------------------------------------------------
A high SHAP importance does not automatically increase the weight of every
row containing that feature: only rows occupying underrepresented regions
receive a positive contribution from it. If all contributing deficits are
zero, every row has weight one and the augmentation distribution is uniform.
Writing $N=|D_{\mathrm{train}}|$ and $m=\lfloor\gamma N\rfloor$, the expected
share of original record $i$ in the augmented training table is
$(1+m\pi_i^{(a)})/(N+m)$. The original copy is retained and the expected
number of additional copies is $m\pi_i^{(a)}$. Thus $w_{\max}$ bounds the
sampling weight, not the realised number of duplicates.

% ----------------------------------------------------------------------
% SECTION 4.2.1: replace the patient-level stay-selection description.
% ----------------------------------------------------------------------
The query first restricted ICU stays to those whose initial care unit was
the Medical Intensive Care Unit (MICU) or the Medical/Surgical Intensive
Care Unit (MICU/SICU). Within this restricted set, stays were ordered by
\texttt{intime} for each \texttt{subject\_id}, and the earliest stay was
selected. The age restriction of 18--89 years was then applied to that
selected stay. The resulting record therefore represents the patient's
first stay in one of the selected care units, not necessarily the patient's
first ICU stay overall. An earlier stay in another ICU does not prevent
inclusion. Conversely, the query does not search for a later age-eligible
stay if the earliest selected-unit stay fails the age restriction. This
selection retains at most one ICU stay per patient in the extracted cohort.

% SECTION 4.2.1: replace the following exclusion sentence.
The SQL step excluded stays outside the selected initial care units,
retained only the earliest stay within those units per patient, and then
excluded patients whose selected stay did not meet the age restriction.
Subsequent preprocessing excluded records with missing values in the
selected analysis variables.

% TABLE 4.1, first operation cell: use the following wording.
% Restrict initial care units; retain the earliest selected-unit stay per
% patient; apply the age restriction; join first-day measurements.
%
% FIGURE 4.1 cohort-extraction box:
% Selected initial care units; earliest stay within these units per patient;
% age restriction; first-day measurements.
%
% APPENDIX A, add before Listing A.1; do not change the executed SQL.
The care-unit restriction precedes the within-patient ranking in Listing A.1.
The age restriction is applied after ranking. Accordingly, the query selects
the earliest stay in the specified initial care units and retains it only
if its age criterion is met; it does not select the first ICU stay overall.

% ----------------------------------------------------------------------
% SECTION 5.3: replace the blanket claim that all transformations fit real train.
% ----------------------------------------------------------------------
Preprocessing was fitted separately for each component without using its
evaluation holdout. Generator transformations used only the applicable
generator-training data, subject to the transformer-reuse policy described
in Section 5.7. The post-hoc detector fitted its preprocessing on its
real/synthetic detector-training fold. Each utility classifier fitted its
preprocessing on its own training sample: synthetic-only, real-only, or
mixed, after any class-balancing resampling. Nearest-neighbour preprocessing
was fitted on the real training reference used for that analysis. These
component-specific transformations were then applied unchanged to their
respective evaluation data.

% ----------------------------------------------------------------------
% TABLE 5.4: replacement caption and final column header.
% Keep all existing coefficient values and rows.
% ----------------------------------------------------------------------
% Caption:
% Ablation variants in the explanation-guided retraining experiment.
% The coefficients define feature-priority scores, not direct row weights.
% All A2--A5 variants combine their selected feature priorities with the same
% distributional region deficits before constructing sampling probabilities.
%
% Last column header: Feature-priority rule / augmentation
% A3 last column: SHAP-only priorities; shared region-deficit row weighting.

% SECTION 5.7: replace the transformer-reuse paragraph.
For non-private DP-CGANS, the transformer fitted for A0 on the original
real training partition was frozen, hash-validated, and reused for A1--A5.
The saved-transformer interface makes this reuse possible without changing
the generative objective and avoids repeated expensive mixture fitting.
Each variant still transformed its own augmented rows, rebuilt the sampler
from those rows, and trained freshly initialised generative networks. Thus
the fitted representation was fixed, but conditional sampling frequencies
were not frozen to their A0 values. This is an engineering and experimental
control choice, not a requirement of differential privacy; DP-CGANS was
used in non-private mode.

CTAB-GAN+ and CTGAN instead refitted their transformers on each variant's
training table. Their ablation effects can therefore include changes in
the fitted representation as well as changes in row exposure. The
transformer policy is not identical across architectures, and the results
compare these implemented pipelines rather than isolating reweighting under
one common fixed-representation policy. A frozen-versus-refitted transformer
control would be needed to separate these effects across all models.

% SECTION 5.7: augmentation size wording.
% Replace "0.6N" and "1.6N" as exact counts with
% m = floor(0.6N) extra rows and N + m total rows, respectively.
% Equal epoch counts do not imply equal optimiser-update counts when the
% training set is enlarged. A1 controls this additional exposure for A2--A5.

% ----------------------------------------------------------------------
% CHAPTER 6: use the corrected local evaluation.tex, or transfer the targeted
% changes listed in audit.md if retaining the version-(6) chapter source.
% Its existing labels are preserved; the added AP equation is unnumbered.
% Key replacement privacy interpretation is repeated here for convenience.
% SECTION 6.3.4: replace "Lower values indicate fewer direct copies...".
% ----------------------------------------------------------------------
A non-zero exact-match rate is not by itself evidence of copying. Matches
can arise by chance, particularly with discrete variables, rounded
measurements, or repeated common profiles; there is no general requirement
that a GAN's exact-match rate be zero. Conversely, zero exact matches only
means that no complete match was observed in the evaluated synthetic
sample. It does not exclude near-copies, partial memorisation, or membership
and attribute inference. The same representation and rounding policy must
be used for both datasets because these choices affect exact equality.

Nearest-neighbour distances provide complementary descriptive evidence,
not a privacy guarantee. They depend on scaling, encoding, dimensionality,
reference-set size, and the density of the underlying population. A
synthetic-to-held-out distance ratio below one warrants inspection but does
not prove memorisation; a ratio near one does not establish safety, and
larger distances may simply indicate poorer fidelity. Aggregate distance
percentiles can conceal exposure of individual patients or small subgroups.
These checks should therefore be distinguished from attack-based privacy
audits under an explicit threat model and from verified differential-privacy
guarantees~\cite{Stadler2022SyntheticData}.

% ----------------------------------------------------------------------
% RESULTS, DISCUSSION, CONCLUSION: terminology replacements only.
% ----------------------------------------------------------------------
% "Jensen--Shannon divergence" -> "Jensen--Shannon distance"
% wherever the term refers to the implemented/reported categorical metric.
% Keep "divergence" when defining the expression inside the square root.
%
% "PR-AUC" -> "AP (PR)" in plot/table headings, and
% "average precision (AP)" at first prose occurrence, then "AP".
% AP is the implemented non-interpolated precision--recall summary.
% Do not change saved pr_auc/utility_pr_auc field names or numerical results.
%
% "Tail-CDF divergence" -> "upper-tail ECDF gap" for Eq. (6.6) and
% Tables 7.3--7.5, with legacy output key documented in Chapter 6.
% Do not relabel fixed-tail mass error or quantile-location error as this gap.
%
% DISCUSSION: optional limitation paragraph.
The ablation labels distinguish the signals used for feature prioritisation,
not completely independent weighting mechanisms. In particular, A3 combines
SHAP-only priorities with the shared region-deficit construction. Differences
between A3 and A2 therefore contrast explanation-based and mismatch-based
feature priorities within that common mechanism. Interpretation across GAN
architectures additionally requires acknowledging the asymmetric transformer
policy: the DP-CGANS representation is fixed across variants, whereas those
of CTAB-GAN+ and CTGAN are refitted.
